% !PRINTED 437
\weberpage{437}
We now consider the sequence of class invariants
\weq{10}{\ldots, (l^{-2}k), (l^{-1}k), (k), (lk), (l^2k), (l^3k), \ldots,}
which can be continued arbitrarily forwards and backwards. In this
sequence, if $l$ belongs to the exponent $\varepsilon$, then $(l^{\nu}k)$ is
identical with $(l^{\nu+\varepsilon}k)$, while all the terms lying between
them differ from it and from one another. If $l$ is two-sided, the
sequence (10) contains only two different terms, each of which is
\gesperrt{rationally expressible} by the other. Otherwise we conclude from (8)
that every term of the sequence (10) is rationally expressible by the
two preceding ones (or also by the two following ones) in the form
\weq{11}{(l^{\nu+1}k) = -(l^{\nu-1}k) + f_l(l^{\nu}k).}
By a repeated application of this formula one arrives at the theorem
\gesperrt{that every term of the sequence (10) can be rationally expressed by
any two consecutive terms of the same sequence.}

\gesperrt{But if $l$ is two-sided, then all the terms of the sequence
(10) are rationally expressible by one of them.}

\begin{center}
§ 121. \gesperrt{Separation of the opposite classes.}
\end{center}

It is now a question of separating the two roots of the quadratic
equation from one another in the case of an $l$ that is not two-sided
as well, which is possible by the adjunction of $\sqrt{D}$.

If the class $l$ is not two-sided, the two class invariants
$(lk)$, $(l^{-1}k)$ are different from one another, and therefore
(§ 118, III.)
\[
v = (lk) = j\left(\frac{\omega}{p}\right)
\]
is a \gesperrt{simple} root of the transformation equation
\[
F_p(u, v) = 0
\]
(§ 119, V.). If we therefore put
\weq{1}{M = \left(\frac{\eta\left(\dfrac{\omega}{p}\right)}{\eta(\omega)}\right)^{24},}
then it follows by § 72 that $M$ is \gesperrt{rationally} expressible by
\[
j(\omega) \quad \text{and} \quad j\left(\frac{\omega}{p}\right),
\]
%CONT>>
